\usepackage[utf8]{inputenc}
\usepackage[OT1]{fontenc}
\usepackage{ae}
\usepackage[brazil]{babel}
\usepackage{a4wide}
\usepackage{comment}
\usepackage[pdftex]{graphicx,color}
\usepackage{graphics}
\usepackage{cite}
\usepackage{longtable}
\usepackage{float}
\usepackage{fancyvrb}
\usepackage{fancyhdr}
\usepackage{setspace}
\usepackage{amsmath}
\usepackage{lscape}
\usepackage{textcase}
\usepackage{anysize}
\usepackage{booktabs}
\usepackage{url}
\usepackage{subfig}
\usepackage{xcolor}

\usepackage{listings}

% Configura o estilo do 'listings' para este documento
\lstset{
    language=Python,                % Define a linguagem (Python)
    basicstyle=\ttfamily\small,     % Define a fonte e o tamanho
    breaklines=true,                % Quebra linhas longas
    showstringspaces=false,         % Não mostra espaços em strings
    captionpos=b,                   % Posição da legenda (bottom)
    frame=tb,                        % Adiciona um quadro simples (top/bottom)
    inputencoding=utf8,
    literate={á}{{\'a}}1 {ã}{{\~a}}1 {â}{{\^a}}1
             {é}{{\'e}}1 {ê}{{\^e}}1
             {í}{{\'i}}1
             {ó}{{\'o}}1 {õ}{{\~o}}1 {ô}{{\^o}}1
             {ú}{{\'u}}1
             {ç}{{\c{c}}}1
             {Á}{{\'A}}1 {Ã}{{\~A}}1 {Â}{{\^A}}1
             {É}{{\'E}}1 {Ê}{{\^E}}1
             {Í}{{\'I}}1
             {Ó}{{\'O}}1 {Õ}{{\~O}}1 {Ô}{{\^O}}1
             {Ú}{{\'U}}1
             {Ç}{{\c{C}}}1,
    keywordstyle=\color{blue},
    commentstyle=\color{gray},
    stringstyle=\color{blue},
}

\usepackage[alf]{abntex2cite}

\marginsize{20mm}{20mm}{20mm}{15mm}


%% Cabeçalhos
\renewcommand{\topfraction}{1}
\renewcommand{\bottomfraction}{1}
\renewcommand{\floatpagefraction}{1}
\renewcommand{\textfraction}{0}
%\renewcommand{\baselinestretch}{2}
\doublespacing %espaçamento duplo
\sloppy

%% Nomes
\floatstyle{plain}  %%% tipos: plain, boxed, ruled
\newfloat{codigo}{tbp}{lop}[section]
\floatname{codigo}{Código}

%%% nome para ser usado no sumário

\newcommand{\listofcodename}{Lista de C\'{o}digos}



% RESUMO ----------------------------------------------------------------------------------------------------------------------------------------------------------------------

\newcommand{\resumo}[1]{
\begin{center} \LARGE \bf Resumo \end{center}

\vskip 4em
\input{#1}

\newpage

}

% ABSTRACT ----------------------------------------------------------------------------------------------------------------------------------------------------------------------

\newcommand{\abstractt}[1]{
\begin{center} \LARGE \bf Abstract \end{center}

\vskip 4em
\input{#1}

\newpage

}

% Sumário -----------
\newcommand{\sumario}{
\renewcommand{\contentsname}{Sum\'{a}rio}
\tableofcontents
\addcontentsline{toc}{chapter}{\listtablename}
\listoftables

\newpage
\addcontentsline{toc}{chapter}{\listfigurename}
\listoffigures
\addcontentsline{toc}{chapter}{\listofcodename}
\listof{codigo}{\listofcodename}  % Lista de Códigos

\clearpage
}
